\section{Nondegeneracy of the transport scalar}
\label{sec:nondegeneracy}

Section~\ref{sec:actual-arcs} extracts a limit $a$ of the bounded scalars
$a_h^Q$ and uses the energy bound to obtain \eqref{transport:bulk} with
this single scalar. Its choice depends only on the bulk subsequence, and
does not depend on any probes used in Section~\ref{sec:observable}. We now
prove the asserted restriction on $a$. The strict area imbalance between
the two Fourier fields supplies the positive quantity; contour estimates
and uniform integrability justify passing it through the transport limit.

\begin{proposition}[Nondegeneracy]
\label{prop:nondegeneracy}
Every scalar obtained from the preceding transport construction satisfies
\[
                 1-2\Rea a>0.
\]
In particular, $a\ne1$.
\end{proposition}

Fix a nonzero $f\in C_c^\infty(\C)$ and work in the full plane.
Write $R=Q(f)-aV(f)$. The consequence of bulk transport needed here is
the averaged pairing limit
\begin{equation}
             \EE\langle R,V(f)\rangle\longrightarrow0.
 \label{energy:pairing}
\end{equation}
Here the inner product is in the random label space
$\ell^2(\{\lambda\})$, with the second factor conjugated. Once
\eqref{energy:pairing} is proved, the exact energy identities for $U$ and
$Q$ will force $1-2\Rea a>0$.

Bulk transport gives convergence in probability on macroscopic rims;
it does not by itself justify this sum over all labels or its expectation.
We therefore replace $V(f)$ by contour integrals. Their squared norm on
small rims is negligible, while area multiplicity and wall-count moments
control their macroscopic sum. These two estimates will make the pairing
uniformly integrable and establish \eqref{energy:pairing}.
For a label $\lambda=(L,C)$ with polygonal rim $P$, put
\[
             Y_\lambda=Y_P=\int_P f(z)\,\dd z;
\]
for an isolated-edge component put $Y_\lambda=0$. All rims are oriented as
in Lemma~\ref{lemma:rim}; norms below use the same random label space.

\subsection{Replacing the tagged tangent field by a contour integral}

We first establish the two estimates
\begin{align}
 \EE\sum_\lambda |Y_\lambda-V_\lambda(f)|^2&\longrightarrow0,
 \label{energy:comparison}\\
 \lim_{b\downarrow0}\limsup_{\eps\to0}
 \EE\sum_{\diam(P)<b}|Y_P|^2&=0.
 \label{energy:small-rims}
\end{align}
Choose a compact box $K_0$ containing $\supp f$ in its interior; enlarge it
by a fixed amount whenever a mesh estimate requires a collar.

For a tagged triangle $t$, its clockwise missing side is the corresponding
oriented rim edge. Comparing $f(z_t)$ with $f$ along that edge costs at
most $C_f\diam(t)^2$. An isolated edge has two opposite oriented edge
integrals, which cancel, and the same estimate applies to its two tags.
Consequently there are nonnegative budgets
\[
 B_\lambda=C_f\!\sum_{\substack{t\text{ tagged by }\lambda\\
                         t\cap\supp f\ne\varnothing}}
                         \diam(t)^2
 \quad\text{such that}\quad
 |Y_\lambda-V_\lambda(f)|\le B_\lambda.
\]
Triangles in a vanishing mesh collar may be included in this sum without
changing the argument. A triangle has at most one tag, so
$S_\eps=\sum_\lambda B_\lambda$ has uniformly bounded moments of every
fixed finite order by Lemma~\ref{lemma:mesh}.

We claim that $M_\eps=\max_\lambda B_\lambda\to0$ in probability.
First fix $b>0$ and tile a neighborhood of $\supp f$ by squares of side
$s\ll b$. If a rim of diameter at least $b$ has a tag in such a square,
then, on fine mesh, its monochromatic path gives an arm from a fixed
enlargement of the square to distance of order $b$. The local
diameter-square sum has $L^2$ norm at most $Cs^2$ in the microscopic upper
limit. Cauchy--Schwarz and the one-arm estimate therefore bound the
expected budget in that square belonging to these rims by
\[
                    C_f s^2(s/b)^{c/2}
\]
in that upper limit. There are $O(s^{-2})$ squares. Thus the total budget
of the rims of diameter at least $b$ vanishes in probability after first
taking $\eps\to0$ and then $s\downarrow0$.

A rim of diameter less than $b$ has all its relevant tags in an enlarged
box of side $O(b)$ on fine mesh. The same is true of an isolated tag pair.
Cover the fixed support neighborhood by $O(b^{-2})$ such boxes. If $D_B$
denotes the diameter-square sum in an enlarged box, then, for every fixed
$p>1$,
\[
 \limsup_{\eps\to0}\PP\!\left(\max_B D_B>u\right)
 \le C_p u^{-p}b^{2p-2},\qquad u>0,
\]
using the local $L^p$ bound $\limsup_{\eps\to0}\|D_B\|_p\le C_pb^2$.
Letting $b\downarrow0$ proves the claim. Finally,
\[
 \sum_\lambda |Y_\lambda-V_\lambda(f)|^2
       \le M_\eps S_\eps.
\]
The right side tends to zero in probability and is uniformly integrable:
$M_\eps\le S_\eps$ and the latter has bounded moments of order greater
than two. This proves \eqref{energy:comparison}.

\subsection{Distinct walls and area multiplicity}

We need two bounds on these contours: a small-diameter estimate for
their expected square sum, and higher moments of their absolute sum
above a fixed diameter cutoff. Both follow by counting distinct walls
and controlling how many of one wall's component interiors cover a point.

Stokes' theorem, applied with either orientation, gives
\begin{equation}
 |Y_P|\le C_f\int_{K_0}\ind_{\operatorname{int}_b(P)}(z)\,\dd^2z,
 \label{energy:stokes}
\end{equation}
where $\operatorname{int}_b(P)$ is the bounded geometric inside of $P$.
Also, $Y_P=0$ unless $P$ meets $K_0$.

For each fixed wall $L$, at most two of its rims enclose a generic point.
Indeed, the bounded components of the untouched triangles have disjoint
disk interiors by Lemma~\ref{lemma:rim}; in addition, there is only one
unbounded component. Its rim can contribute one further bounded inside.
Polygonal boundaries have area zero, so this pointwise assertion is
sufficient for all area integrals below. It is an area-multiplicity
statement, and does not bound the number of rims associated with a wall.

We record the wall-count estimate used at both small and fixed scales.
At any scale $t>0$, the expected number of rims of diameter in $[t,2t]$
enclosing a specified generic point is bounded uniformly in $t$ and
$\eps$. Such a rim lies within distance $2t$ of the point. On mesh fine
relative to $t$, its wall visits the triangles facing the rim edges and
therefore has a spanning arc in one of finitely many fixed-ratio annuli
centered at a grid of mesh a small fixed multiple of $t$. Nested primal
boundary cycles with room between them may be chosen as in
Section~\ref{sec:probability}. For each such annulus, let $N$ count the
proper spanning wall strands obtained by cutting at all boundary
contacts and stopping at the boundary edges. Use the bank extraction
from the proof of Lemma~\ref{lemma:mixed-four}: every strand has an
adjacent black bank, each gap between consecutive strands accounts for
at most two such banks, and banks selected from distinct gaps give
vertex-disjoint black crossings after proper trimming. This includes
their boundary endpoints, by the stopped-edge argument in that proof:
each black boundary site belongs to only one gap.
In particular, $N\ge2k$ supplies $k$ crossings with
disjoint site witnesses. On the crossing-good geometry event $G$, a
fixed buffered subannulus has quenched white-circuit probability at
least $\delta>0$, supplied by finitely many fixed-aspect rectangle tests.
Such a circuit excludes a black spanning crossing. BK, applied
conditionally on the point geometry, therefore gives
\[
            \Peta(N\ge2k)\le(1-\delta)^k\qquad\text{on }G.
\]
Consequently the spanning-strand counts have bounded quenched moments
of every fixed order on $G$.

Here and below the exceptional geometry is controlled quantitatively.
For $t/\eps$ large, the probability that either the needed mesh or crossing
tests fail is at most $C_\gamma(t/\eps)^{-\gamma}$, for any prescribed
finite $\gamma$, by Proposition~\ref{prop:probability-input} and
Lemma~\ref{lemma:mesh}. Without these tests, count distinct walls by the
number $J$ of incident triangles in the relevant $O(t)$ region. Every
counted wall has a tagged incident triangle there, and the moments of $J$
grow at most polynomially in $1+t/\eps$. For example, if
\[
 \|J\|_{2m}\le C(1+t/\eps)^{k_m},
\]
then
\[
 \EE[J^m\ind_{\mathrm{bad}}]
 \le \|J\|_{2m}^{m}\PP(\mathrm{bad})^{1/2}
 \le C(1+t/\eps)^{mk_m}(t/\eps)^{-\gamma/2}.
\]
Choose $\gamma>2mk_m$. For bounded $t/\eps$, the incidence bound alone
suffices. The two-per-wall area multiplicity gives the asserted uniform
expected rim count.

Since $|Y_P|\le C_ft^2$ for $\diam(P)\in[t,2t]$, this count and
\eqref{energy:stokes} imply
\[
 \EE\sum_{t\le\diam(P)<2t}|Y_P|^2
 \le C_ft^2\int_{K_0}
       \EE\sum_{t\le\diam(P)<2t}
            \ind_{\operatorname{int}_b(P)}(z)\,\dd^2z
 \le C_f t^2.
\]
Dyadic summation over scales below $b$ proves
\eqref{energy:small-rims}.

For fixed $b>0$, let $N_{\eps,b}$ count the distinct walls represented by
rims with $\diam(P)\ge b$ and $Y_P\ne0$. Cover $K_0$ by finitely many grid
balls whose radii are a sufficiently small fixed fraction of $b$. A rim
of diameter at least $b$ meeting one of these balls must leave a fixed
enlargement of it, so its wall has a spanning arc in a corresponding
fixed-ratio annulus. The preceding spanning-arc and bad-geometry estimates,
now at fixed scale, prove
\[
          \sup_{\eps\text{ small}}\EE N_{\eps,b}^{m}<\infty
          \qquad (b>0,\ m<\infty).
\]
The constants may depend on $b$ and $m$. Put $M_f=C_f\Area(K_0)$.
Applying area multiplicity before summing labels gives the precise bounds
\begin{align}
 \sum_{\diam(P)\ge b}|Y_P|&\le 2M_fN_{\eps,b},\notag\\
 \sum_{\diam(P)\ge b}|Y_P|^2&\le 2M_f^2N_{\eps,b}.
 \label{energy:macro-multiplicity}
\end{align}
In the second line we used $|Y_P|\le M_f$ individually. Thus the absolute
sum, and the square sum, have all fixed higher moments at fixed $b$,
without a count of all rims belonging to one wall.

\subsection{Uniform integrability and the positive area identity}

Put $A_\eps=\|R\|_{\ell^2}$.
The energy identity gives $\sup_\eps\|A_\eps\|_2<\infty$.
At fixed $b$, split $f$ into finitely many smooth tests supported on disks
small compared with $b$. Every relevant rim leaves the fourfold disk of
each such test that it meets. Hence \eqref{transport:bulk} gives
\[
 \delta_{\eps,b}:=\max_{\substack{\diam(P)\ge b\\Y_P\ne0}}
                           |R_P|\longrightarrow0
 \quad\text{in probability}.
\]
By \eqref{energy:macro-multiplicity},
\[
 \left|\sum_{\diam(P)\ge b}R_P\overline{Y_P}\right|
 \le\delta_{\eps,b}\sum_{\diam(P)\ge b}|Y_P|
 \longrightarrow0
 \quad\text{in probability}.
\]
For the required passage to expectation, define
$B_{\eps,b}=(\sum_{\diam(P)\ge b}|Y_P|^2)^{1/2}$.
For any fixed $q>2$ the wall moments give
$\sup_\eps\|B_{\eps,b}\|_q<\infty$. With
$p=2q/(q+2)>1$, Cauchy--Schwarz in labels and H\"older in probability give
\begin{equation}
 \left\|\sum_{\diam(P)\ge b}R_P\overline{Y_P}\right\|_p
 \le\|A_\eps B_{\eps,b}\|_p
 \le\|A_\eps\|_2\|B_{\eps,b}\|_q\le C_{b,f}.
 \label{energy:macro-ui}
\end{equation}
This is uniform integrability, so the macroscopic pairing tends to zero
in $L^1$. The two errors left in replacing this pairing by the full
pairing with $V(f)$ satisfy
\begin{align*}
 \EE|\langle R,V(f)-Y\rangle|
 &\le (\EE\|R\|^2)^{1/2}
            (\EE\|V(f)-Y\|^2)^{1/2},\\
 \EE\left|\sum_{\diam(P)<b}R_P\overline{Y_P}\right|
 &\le (\EE\|R\|^2)^{1/2}
       \left(\EE\sum_{\diam(P)<b}|Y_P|^2\right)^{1/2}.
\end{align*}
The first vanishes as $\eps\to0$ by \eqref{energy:comparison}; the second
vanishes after the microscopic upper limit and then $b\downarrow0$ by
\eqref{energy:small-rims}. This proves the pairing limit
\eqref{energy:pairing}.

\begin{proof}[Proof of Proposition~\ref{prop:nondegeneracy}]
Since $V=U+Q$, the cross terms in the following real part cancel exactly.
By \eqref{fields:energy} and the triangle area identity,
\begin{align*}
 \EE\Re\sum_\lambda(V_\lambda(f)-2Q_\lambda(f))
                                      \overline{V_\lambda(f)}
 &=\EE\sum_\lambda\bigl(|U_\lambda(f)|^2-|Q_\lambda(f)|^2\bigr)\\
 &=\sqrt3\,\EE\sum_t\Area(t)|f(z_t)|^2
 \longrightarrow\sqrt3\int_\C|f(z)|^2\,\dd^2z>0.
\end{align*}
The final approximation follows from fine mesh and the diameter moments.
On the other hand, \eqref{energy:pairing} writes the left side as
\[
              (1-2\Rea a)\,\EE\|V(f)\|^2+o(1).
\]
The norm is nonnegative and has bounded expectation. A nonpositive
coefficient is incompatible with the strictly positive limit. Therefore
$1-2\Rea a>0$, as claimed.
\end{proof}
